\subsection*{Cross-cohort transfer of the recoverability verdict}
\label{sec:transfer}
The certificate is computed from a \emph{fitted} dipolar subspace $\bm M_s$ (the rank-3 PCA of one
cohort's segment-$s$ samples). A natural worry is that the verdict is an artefact of the fitting
cohort. It is not: because $\eta_{s,\ell}(S)=\|\bm e_\ell^\top\bm M_s(\bm I-\bm P_{\mathrm{obs}})\|_2$
depends on $\bm M_s$ \emph{only through its column space}, it is invariant to any rotation of the
$3$-dimensional dipole basis, and drifts across cohorts only as fast as the two cohorts' dipolar
\emph{subspaces} separate.

\begin{theoremc}[Cross-cohort transfer bound]\label{thm:transfer}
Let cohorts $A,B$ have segment-$s$ dipolar bases $\bm M_s^A,\bm M_s^B\in\mathbb R^{12\times3}$ with
orthonormal columns, and let $\theta^\star$ be the largest principal angle between
$\operatorname{col}(\bm M_s^A)$ and $\operatorname{col}(\bm M_s^B)$. For every target lead $\ell$
and observed set $S$,
\[
\big|\eta_{s,\ell}^A(S)-\eta_{s,\ell}^B(S)\big|
\;\le\; 2\sin(\theta^\star/2)\,\big(1+2\,\kappa_{s}(S)\big),
\qquad \kappa_s(S)=\|\bm M_{s,S}^{+}\|_2=\|\bm M_s\bm M_{s,S}^{+}\|_2 ,
\]
the observed-set conditioning of Prop.~\ref{thm:perlead}. In particular $\theta^\star{=}0$ gives
$\eta^A_{s,\ell}(S)=\eta^B_{s,\ell}(S)$ exactly.
\end{theoremc}

\noindent\emph{Proof sketch.} $\eta$ is unchanged if $\bm M_s$ is right-multiplied by any
$\bm R\in O(3)$ (it rotates $\bm P_{\mathrm{obs}}$ and $\bm e_\ell^\top\bm M_s$ together), so align the
bases by the optimal $\bm R$, giving $\|\bm M_s^A-\bm M_s^B\|_2=2\sin(\theta^\star/2)=:\delta$. With
$\bm m_\ell=\bm M_s^\top\bm e_\ell$ ($\|\bm m_\ell\|\le1$) and $\bm P_S$ the row-space projector,
$\big|\,\|(\bm I{-}\bm P_S^A)\bm m_\ell^A\|-\|(\bm I{-}\bm P_S^B)\bm m_\ell^B\|\,\big|
\le\|\bm m_\ell^A-\bm m_\ell^B\|+\|\bm P_S^A-\bm P_S^B\|_2\,\|\bm m_\ell^B\|
\le\delta+2\delta\kappa_s(S)$, using the projector perturbation
$\|\Delta\bm P_S\|_2\le2\|\Delta\bm M_{s,S}\|_2\|\bm M_{s,S}^{+}\|_2$ (Stewart--Wedin) and
$\|\bm m_\ell^B\|\le1$.\hfill$\square$

\smallskip
\noindent\textbf{What this buys, and its honest limits.} The bound is a $\sin\Theta$ perturbation
(Davis--Kahan~\cite{davis1970rotation}) that makes a \emph{single} extra cohort load-bearing: the
transferability of the verdict is governed by a \emph{computable} quantity---the principal angle
between the cohorts' dipolar subspaces---so a large angle \emph{flags} non-transfer before any
target-cohort label is seen. The constant $1+2\kappa_s$ is worst-case and the bound is loose in
absolute terms (for a well-conditioned observed set it can exceed the trivial $\eta\le1$); its value
is as a \emph{pre-deployment transfer-risk flag}: a small principal angle certifies a small
verdict drift, a large one refuses to. On PTB-XL$\to$Chapman (matched sinus, limb-6 $\to$ V2) the two
regimes are visible. The QRS dipolar subspace is nearly shared
($\theta^\star{=}\TransferAngleQRS^\circ$) and the measured drift is correspondingly tight
($\TransferDriftQRS$, realized slope $\TransferLipQRS$)---the verdict transfers. The ST subspace is
near-orthogonal ($\theta^\star{=}\TransferAngleST^\circ$, CI up to $89.6^\circ$): the bound turns
vacuous and the drift is genuinely uncertain---the target-cohort $\eta$ bootstrap CI is
$[0.02,0.71]$, spanning most of the admissible range, so the realized slope is only pinned to
$[\TransferSTLipLo,\TransferSTLipHi]$ and we claim
\emph{no} precise ST drift, only that the large angle flags the ST verdict as non-transferable
(\texttt{experiments/transfer\_bound.py}). We therefore do not read a stable cross-segment Lipschitz
constant from these two points; the load-bearing claim is the rigorous bound plus the angle as a
computable transfer-risk diagnostic, corroborated by the tight QRS point. This is the cross-cohort
analogue of the finite-sample $\sin\Theta$ term in Thm.~\ref{thm:floor}(iii): both control how a
fitted subspace perturbs the verdict.
